
\subsubsection{3.1 Overview}

The methodology comprises three components: benchmark selection for ground-truth labels, controlled variables for fair comparison, and mechanism verification to separate implementation correctness from discrimination effectiveness.

\subsubsection{3.2 Benchmark}

TruthfulQA mc1 (multiple-choice, single correct answer) serves as the evaluation benchmark. The multiple-choice format provides ground-truth labels without human annotation: the model's selected answer is either correct or incorrect, enabling AUROC computation where incorrect answers constitute hallucinations.

\subsubsection{3.3 Model}

Llama-3-8B-Instruct serves as the evaluation model. This model is representative of the decoder-only instruction-tuned class, with open weights enabling reproducibility.

\subsubsection{3.4 UQ Methods}

Three UQ methods are evaluated:

\textbf{Max Probability (Token-Level):} Uncertainty score $u = 1 - \max_c P(c)$ where $c \in \{A, B, C, D\}$.

\textbf{Choice Entropy (Token-Level):} Uncertainty score $u = H(P) = -\sum_c P(c) \log P(c)$ over answer choices.

\textbf{Semantic Entropy (Semantic-Level):} Uncertainty score $u = H(P_{\text{cluster}})$ computed over NLI-clustered responses from $N=5$ samples at temperature 0.7. The NLI model is facebook/bart-large-mnli with entailment threshold 0.7.

\subsubsection{3.5 Mechanism Verification}

To distinguish implementation error from format unsuitability, mechanism verification checks whether semantic clustering functions correctly independent of discrimination:

\begin{itemize}
\item \textbf{Cluster count:} Average clusters per question should be less than the number of samples if clustering occurs
\item \textbf{Entropy variance:} Standard deviation of entropy should be positive if the signal varies
\end{itemize}

If mechanism verification passes but discrimination fails, the method is correctly implemented but unsuited to the task format.

